\section{Target Application: Distributed VLIW Computing}
\label{sec:target}

\subsection{VLIW Cores and Compile-Time Scheduling}

VLIW processors expose multiple functional units to the compiler, which schedules
operations explicitly in long instruction words~\cite{ellis1986}. This design
philosophy shifts control complexity from hardware to software, resulting in
simple, area-efficient cores well suited for replication in many-core arrays.
Signal processing, scientific computing, and neural network inference kernels are
representative workloads.

In a distributed VLIW many-core system, each core executes a portion of a
data-flow graph, passing intermediate results to neighboring cores through the
NoC. The communication graph is typically known at compile time, enabling the
compiler to:
\begin{enumerate}
  \item Assign each communication a source route through the NoC;
  \item Schedule packet injections to avoid simultaneous contention on the same
        egress port;
  \item Guarantee bounded communication latency for critical paths.
\end{enumerate}

This compile-time visibility makes HyNoC's hybrid model particularly effective:
the circuit-establishment latency is paid once per packet and is deterministic
(bounded by route length and PRRA grant latency), and the wormhole payload
transfer proceeds without further arbitration overhead.

\subsection{Integration with the Local Interface}

Each VLIW core connects to the NoC through a local interface, which presents a
simple FIFO-based write/read API. The core writes routing flits followed by
payload flits into the local interface's ingress FIFO; the local interface manages
all handshaking with the router. On the receive side, incoming flits are deposited
in the egress FIFO and read by the core at its own pace. The dual-FIFO structure
fully decouples the core's clock domain from the router domain.

The \texttt{hynoc\_stream\_writer} and \texttt{hynoc\_stream\_reader} modules are
simulation utilities that generate and check random packet streams, respectively.
They serve as traffic generators in the testbench environment, verifying the
correctness of the NoC under various load conditions.

\subsection{Multicast for Broadcast Operations}

VLIW many-core applications frequently require broadcast operations: distributing
a coefficient matrix, a lookup table, or a control token to multiple cores
simultaneously. HyNoC's multicast routing protocol supports this natively: a
single packet with a multicast routing flit is replicated at each router toward
all targeted egress ports, without requiring the sender to issue multiple unicast
packets. This reduces network load and sender overhead for one-to-many
communication patterns common in data-parallel VLIW workloads.
